% Generated draft: requires mathematical and source review.
\documentclass[11pt]{article}
\usepackage[margin=25mm]{geometry}
\usepackage{fontspec}
\setmainfont{Latin Modern Roman}
\usepackage{amsmath,amssymb,mathtools,unicode-math}
\setmathfont{Latin Modern Math}
\usepackage{enumitem,needspace,xurl,hyperref}
\setlength{\parindent}{0pt}
\setlength{\parskip}{5pt}
\setlength{\emergencystretch}{3em}
\newcommand{\N}{\mathbb{N}}
\newcommand{\Z}{\mathbb{Z}}
\newcommand{\Q}{\mathbb{Q}}
\newcommand{\R}{\mathbb{R}}
\newcommand{\C}{\mathbb{C}}
\newcommand{\F}{\mathbb{F}}
\newcommand{\A}{\mathbb{A}}
\newcommand{\PP}{\mathbb{P}}
\newcommand{\E}{\mathbb{E}}
\newcommand{\eps}{\varepsilon}
\newcommand{\dd}{\,\mathrm{d}}
\begin{document}
\section*{UnsolvedMath ID: 20000541}
\textbf{Source problem code:} catalogue entry
\subsection{Definitions and mathematical statement}
Let \(n\ge2\) and let \(\mathcal T_1,\ldots,\mathcal T_n\) be finite collections of doubly infinite tubes in \(\R^n\). For a tube \(T\), let \(w(T)\) be a unit vector along its axis, and \(\chi_T\) its indicator. The supplied reconstruction of (0.20) is
\[
 K(\mathcal T_1,\ldots,\mathcal T_n)=
 \int_{\R^n}\left[\sum_{T_1\in\mathcal T_1,\ldots,T_n\in\mathcal T_n}
 |\det(w(T_1),\ldots,w(T_n))|\prod_{i=1}^n\chi_{T_i}(x)\right]^{1/(n-1)}dx
 \le C_n\prod_{i=1}^n(\#\mathcal T_i)^{1/(n-1)}.
\]
The determinant belongs inside the tuple sum and is unsigned; the extracted product of sums with unbound tube symbols is not a precise alternative formula. Iliopoulou asks whether one can take \(C_n=1\) for ``width one'' tubes. Neither the cross-section shape nor whether width means radius or total diameter is defined, so the numerical question depends on normalization.

The supplied normalization theorem states that for planar strips of common total width \(w>0\), the exact best constant is \(C_2(w)=w^2\). Hence total-width-one strips have constant one, whereas radius-one strips, of total width two, have constant four. The report identifies the latter with a Carbery--Valdimarsson convention; it refutes constant one under that interpretation only. More generally, for a fixed cross-section family scaled by radius factor \(r>0\), it gives \(C_n(r)=r^n C_n(1)\).

Its fixed-direction subclass theorem takes a single direction \(v_i\) per family. A cylinder with direction \(v_i\) is the inverse image of a measurable set in \(v_i^\perp\) under orthogonal projection, and its transverse mass is the \((n-1)\)-dimensional measure of that set. If every tube in family \(i\) has transverse mass at most \(m_i\), the report asserts
\[
 K(\mathcal T_1,\ldots,\mathcal T_n)
 \le\prod_{i=1}^n(m_i\#\mathcal T_i)^{1/(n-1)}.
\]
It attributes the proof to an affine Loomis--Whitney argument and reports sharpness on volume-one cubical grids. It also claims that finite dense lattices of round tubes approach the corresponding mass factor. This concerns fixed direction within each family; the transverse-mass-one problem with varying directions for \(n\ge3\) remains unresolved in the supplied report. No normalization-independent resolution of the original constant-one question is asserted.
\subsection{Short English statement}
Can the determinant-weighted multilinear tube inequality have constant one? The supplied report shows that the answer depends on what width one means, gives exact planar scaling and a sharp fixed-direction subclass, and leaves higher-dimensional varying-direction tubes with normalized transverse mass unresolved.
\subsection{Sources}
\begin{itemize}
\item Local UnsolvedMath registry, permanent ID \(20000541\), source code \texttt{catalogue entry}; AIM analysis corpus, \texttt{aim-analysis-notes.json}, notes[72]. Attribution: Ulam AI / UnsolvedMath Contributors, CC BY 4.0.
\item AIM \emph{Carleson theorems and multilinear operators}, May 18--22, 2015, Question 14 attributed to Marina Iliopoulou, equation (0.20): \url{https://aimath.org/pastworkshops/multilinopsproblems.pdf}. The supplied verification refers to official TeX and doubly infinite tubes but stops before its announced full formula. Original remarks and literature fields are empty.
\item The research summary mentions Carbery--Valdimarsson's radius-one convention and an affine Loomis--Whitney argument without complete bibliographic references.
\item Registry dated August 14, 2026: statement exact, status \texttt{partially\_solved}, classification \texttt{PARTIAL-PROGRESS}; research attempt 2 records problem status unknown, partial result, rigor proved, novelty confidence low. The normalization and fixed-direction results are machine-generated research aids, not peer reviewed, requiring independent verification before citation.
\end{itemize}
\end{document}
